\chapter{Monoidal and enriched cats}\label{ch:monoidal}

\begin{assumption}\label{ass:ch9}
In \cref{sec:feline-tensor,sec:graded,sec:self-enrichment}, $\cC$ is a category
with initial and terminal objects $0$ and $1$. In \cref{sec:feline-tensor,sec:graded}
it carries a monoidal structure $(\otimes,I,a,l,r)$, with associator
$a_{X,Y,Z}\colon(X\otimes Y)\otimes Z\to X\otimes(Y\otimes Z)$ and unitors
$l_X\colon I\otimes X\to X$, $r_X\colon X\otimes I\to X$. The hypotheses of
\cref{sec:additive} are stated there.
\end{assumption}

\section{The feline tensor product}\label{sec:feline-tensor}

\begin{definition}\label{def:feline-tensor}
The \emph{feline tensor product} of cats $K=\langle L_\alpha\rangle_{\alpha\in A}$
and $K'=\langle L'_\beta\rangle_{\beta\in B}$ (normalized without loss of
generality) is
\[
K\otimes K'\defeq\bigl\langle L_\alpha\otimes L'_\beta\bigr\rangle_{(\alpha,\beta)\in A\times B},
\]
and on morphisms $(\varphi,F)\otimes(\psi,G)\defeq(\varphi\times\psi,(F_\alpha\otimes G_\beta)_{(\alpha,\beta)})$.
The \emph{feline unit} is $\mathbf I\defeq\iota(I)$. The associator
$\mathbf a_{K,K',K''}$ has gait $((\alpha,\beta),\gamma)\mapsto(\alpha,(\beta,\gamma))$
and components $a_{L_\alpha,L'_\beta,L''_\gamma}$; the unitors $\mathbf l_K$ and
$\mathbf r_K$ have gaits $(*,\alpha)\mapsto\alpha$ and $(\alpha,*)\mapsto\alpha$
and components $l_{L_\alpha}$ and $r_{L_\alpha}$.
\end{definition}

\begin{theorem}\label{thm:monoidal}
$(\Fel(\cC),\otimes,\mathbf I,\mathbf a,\mathbf l,\mathbf r)$ is a monoidal
category. If $\cC$ is braided (resp.\ symmetric) with braiding $c$, then
$\Fel(\cC)$ is braided (resp.\ symmetric) with braiding of gait
$(\alpha,\beta)\mapsto(\beta,\alpha)$ and components $c_{L_\alpha,L'_\beta}$.
The leg functor is a strict monoidal functor
$(\Fel(\cC),\otimes,\mathbf I)\to(\Set,\times,1)$, and $\iota$ is a strong
monoidal functor.
\end{theorem}

\begin{proof}
Functoriality of $\otimes$ follows from that of $\times$ on $\Set$ and of
$\otimes$ on $\cC$. The associator and unitors are isomorphisms by
\cref{prop:iso}, and natural because their gaits are natural in $\Set$ and
their components are natural in $\cC$. Every coherence axiom (the pentagon and
triangle axioms, and in the braided case the hexagon axioms and, in the
symmetric case, $c_{Y,X}c_{X,Y}=\id$) is an equation between morphisms of
cats; by \cref{prop:automatic} such an equation holds if and only if it holds
on gaits and on each leg component. On gaits it is the corresponding axiom for
the cartesian monoidal structure of $\Set$, and on leg components it is the
corresponding axiom in $\cC$. The leg functor sends $K\otimes K'$ to $A\times B$,
$\mathbf I$ to $1$, and the structural isomorphisms to those of $\Set$, so it is
strict monoidal. Finally $\iota X\otimes\iota Y=\langle X\otimes Y\rangle_{1\times1}$,
which is isomorphic to $\iota(X\otimes Y)$ by the isomorphism with gait
$1\times1\to1$ and identity component; these isomorphisms, together with
$\mathbf I=\iota I$, make $\iota$ strong monoidal, the coherence conditions
again reducing to identities on gaits and components.
\end{proof}

\begin{proposition}\label{prop:tensor-distributes}
For every cat $K$, the functors $K\otimes-$ and $-\otimes K$ preserve small
coproducts. In particular $K\otimes\Legless\iso\Legless\iso\Legless\otimes K$.
This holds whether or not $\otimes$ preserves coproducts in $\cC$.
\end{proposition}

\begin{proof}
For a family $(K_i)$ with leg sets $B_i$, the canonical morphism
$\coprod_i(K\otimes K_i)\to K\otimes\coprod_iK_i$ has gait the canonical bijection
$\coprod_i(A\times B_i)\to A\times\coprod_iB_i$ and identity components, so it
is an isomorphism by \cref{prop:iso}.
\end{proof}

\begin{remark}\label{rem:tensor-chloe}
If the tensor product of two connected objects of $\cC$ is connected, then the
feline tensor product of two cats satisfying the Chloe Condition satisfies it,
since the legs of $K\otimes K'$ are the objects $L_\alpha\otimes L'_\beta$. This holds
for $(\Set,\times)$ and $(\Top,\times)$ but not for $(\Graph,\times)$; the
consequences for Gideon Decompositions are drawn in \cref{sec:cmg-tensor}.
\end{remark}


\begin{proposition}\label{prop:body-monoidal}
Suppose $\cC$ has small coproducts. For cats $K$ and $K'$ let
$\mu_{K,K'}\colon\body(K\otimes K')\to\body K\otimes\body K'$ be the morphism whose
restriction to the summand $L_\alpha\otimes L'_\beta$ is $\inj_\alpha\otimes\inj_\beta$.
Then $\mu_{K,K'}$ is an isomorphism for all $K,K'$ if and only if $\otimes$
preserves small coproducts in each variable separately. In that case
$\body$, with the inverses of the $\mu_{K,K'}$ and the canonical isomorphism
$\body\mathbf I=\coprod_1I\iso I$, is a strong monoidal functor.
\end{proposition}

\begin{proof}
If $\otimes$ preserves coproducts in each variable, $\mu_{K,K'}$ is the
composite of the canonical isomorphisms
$\coprod_{\alpha,\beta}L_\alpha\otimes L'_\beta\iso\coprod_\alpha(L_\alpha\otimes\coprod_\beta L'_\beta)\iso(\coprod_\alpha L_\alpha)\otimes\coprod_\beta L'_\beta$.
Conversely, taking $K=\iota X$ shows that the canonical map
$\coprod_\beta X\otimes L'_\beta\to X\otimes\coprod_\beta L'_\beta$ is invertible for
every family $(L'_\beta)$, and similarly in the other variable. The coherence
conditions for a strong monoidal functor hold because both sides are induced
by the coproduct injections.
\end{proof}

\begin{remark}\label{rem:naive-tensor}
One might attempt to define the tensor product of cats componentwise, as
$(H\otimes H',T\otimes T',(L_\alpha\otimes L'_\beta))$. This is a cat for all
$K,K'$ if and only if $0\otimes0$ is initial and $1\otimes1$ is terminal, since
any initial object is isomorphic to $0$ and any terminal object to $1$. Neither
condition holds in general. In $\Set^{\op}$ with its cartesian monoidal
structure, the product of two copies of the initial object (a singleton) is a
two-element set, which is not initial in $\Set^{\op}$. In $\Set$ with the
cocartesian monoidal structure $\sqcup$, the tensor $1\sqcup1$ of two copies of
the terminal object is not terminal. This is why \cref{def:feline-tensor}
renormalizes the head and the tail.
\end{remark}

\section{Monoidal cats and graded monoids}\label{sec:graded}

\begin{definition}\label{def:monoidal-cat}
A \emph{monoidal $\cC$-cat} is a monoid in the monoidal category
$(\Fel(\cC),\otimes,\mathbf I)$: a cat $K$ with morphisms
$\mu\colon K\otimes K\to K$ and $\eta\colon\mathbf I\to K$ satisfying the
associativity and unit axioms. A morphism of monoidal cats is a monoid
morphism. The category of monoidal $\cC$-cats is denoted $\Mon(\Fel(\cC))$.
\end{definition}

\begin{definition}\label{def:graded-monoid}
Let $M$ be a monoid in $\Set$ with unit $e$. An \emph{$M$-graded monoid} in
$\cC$ is a family $(X_m)_{m\in M}$ of objects with morphisms
$\mu_{m,n}\colon X_m\otimes X_n\to X_{mn}$ and $\eta\colon I\to X_e$ such that
for all $m,n,p\in M$
\begin{gather*}
\mu_{mn,p}\circ(\mu_{m,n}\otimes\id)=\mu_{m,np}\circ(\id\otimes\mu_{n,p})\circ a,\\
\mu_{e,m}\circ(\eta\otimes\id)=l,\qquad\mu_{m,e}\circ(\id\otimes\eta)=r.
\end{gather*}
A morphism $(M,X)\to(N,Y)$ of graded monoids is a monoid homomorphism
$f\colon M\to N$ together with morphisms $g_m\colon X_m\to Y_{f(m)}$ commuting with
the multiplications and units. We write $\GrMon(\cC)$ for the resulting category.
\end{definition}

\begin{theorem}\label{thm:graded-monoids}
There is an equivalence $\Mon(\Fel(\cC))\eqv\GrMon(\cC)$, which is an
isomorphism on normalized cats. Under it, the underlying set of legs of a
monoidal cat, with the multiplication given by the gait of $\mu$, is the
grading monoid.
\end{theorem}

\begin{proof}
Let $(K,\mu,\eta)$ be a monoidal cat with $K=\langle L_\alpha\rangle_A$. The
morphism $\mu$ has a gait $m\colon A\times A\to A$ and components
$\mu_{\alpha,\beta}\colon L_\alpha\otimes L_\beta\to L_{m(\alpha,\beta)}$; the
morphism $\eta$ has a gait $1\to A$, picking out $e\in A$, and a component
$\eta\colon I\to L_e$. By \cref{prop:automatic}, the associativity and unit
axioms of the monoid $K$ hold if and only if they hold on gaits and on
components. Since $\legs$ is strict monoidal (\cref{thm:monoidal}), the axioms
on gaits say that $(A,m,e)$ is a monoid in $\Set$, and the axioms on components
are exactly those of \cref{def:graded-monoid}. Similarly, a morphism of
monoidal cats is a morphism of cats whose gait is a monoid homomorphism and
whose components satisfy the compatibilities of \cref{def:graded-monoid}.
\end{proof}

\begin{example}\label{ex:graded}
\begin{enumerate}
\item For $(\cC,\otimes)=(\Ab,\otimes_{\mathbb Z})$, an $M$-graded monoid is
  a family of abelian groups $(X_m)$ with bilinear associative unital
  multiplications $X_m\times X_n\to X_{mn}$; taking the direct sum, this is the
  same as an $M$-graded ring. Thus monoidal $\Ab$-cats with leg monoid $M$ are
  $M$-graded rings, and monoidal $\Ab$-cats with leg monoid $\mathbb N$ are
  $\mathbb N$-graded rings.
\item For $(\cC,\otimes)=(\Set,\times)$, an $M$-graded monoid in $\Set$ is the
  same as a monoid $N=\coprod_mX_m$ together with a monoid homomorphism $N\to M$
  (sending $X_m$ to $m$). Thus $\Mon(\Fel(\Set))$ is equivalent to the arrow
  category of the category of monoids, in accordance with $\Fel(\Set)\eqv\Set^{\to}$
  (\cref{cor:Fel-Set}).
\item For $\cC=\one$, a monoidal $\one$-cat is a monoid in $\Set$.
\end{enumerate}
\end{example}

\begin{remark}\label{rem:graded-honest}
\Cref{thm:graded-monoids} identifies monoidal cats with graded monoids, a
classical notion. When $\cC$ is symmetric, commutative monoids in $\Fel(\cC)$
correspond to graded monoids with commutative grading monoid and commutative
multiplication in the sense determined by the symmetry of $\cC$. Any signs
in this commutativity come from the symmetry of $\cC$ (for instance, the
Koszul signs of chain complexes involve the homological degree of the legs),
not from the grading monoid: a commutative monoid in $\Fel(\Ab)$ with leg
monoid $\mathbb N$ is a commutative, not a graded-commutative, $\mathbb N$-graded
ring.
\end{remark}

\section{Self-enrichment}\label{sec:self-enrichment}

The category $\Fel(\Set)$ has finite products (\cref{thm:limits}), and so is a
cartesian monoidal category. The following theorem shows that every category
of cats is enriched in $\Set$-cats.

\begin{theorem}\label{thm:self-enrichment}
For $\cC$-cats $K=\langle L_\alpha\rangle_A$ and $K'=\langle L'_\beta\rangle_B$, let
\[
[K,K']\defeq\Bigl\langle\prod_{\alpha\in A}\cC(L_\alpha,L'_{\varphi(\alpha)})\Bigr\rangle_{\varphi\in\Set(A,B)}\in\Fel(\Set).
\]
Let $\mathrm c\colon[K',K'']\times[K,K']\to[K,K'']$ be the morphism of
$\Set$-cats with gait $(\psi,\varphi)\mapsto\psi\varphi$ and components
$(G,F)\mapsto(G_{\varphi(\alpha)}F_\alpha)_\alpha$, and let
$\mathrm u\colon\Caudal\to[K,K]$ have gait $*\mapsto\id_A$ and component
$*\mapsto(\id_{L_\alpha})_\alpha$. These data make $\Fel(\cC)$ a category
enriched in $(\Fel(\Set),\times,\Caudal)$, whose underlying ordinary category,
computed with $\points=\Fel(\Set)(\Caudal,-)$, is $\Fel(\cC)$.
\end{theorem}

\begin{proof}
By \cref{cor:limit-formulas}, $[K',K'']\times[K,K']$ has legs indexed by pairs
$(\psi,\varphi)$, with leg the product of the corresponding legs, so
$\mathrm c$ is well typed. The associativity and unit axioms of an enriched
category are equations between morphisms of $\Set$-cats, and by
\cref{prop:automatic} they reduce to equations on gaits, which are
associativity and unitality of composition of functions, and on components,
which are associativity and unitality of composition in $\cC$. Finally,
$\points[K,K']=\coprod_\varphi\prod_\alpha\cC(L_\alpha,L'_{\varphi(\alpha)})\iso\Fel(\cC)(K,K')$
by \cref{cor:Fel-hom}, and the composition induced by $\mathrm c$ on points is
that of $\Fel(\cC)$.
\end{proof}

\begin{remark}\label{rem:self-enrichment-honest}
Under $\Fel(\Set)\eqv\Set^{\to}$, a category enriched in $(\Set^{\to},\times)$
amounts to a pair of ordinary categories with the same objects and an
identity-on-objects functor between them. The enrichment of
\cref{thm:self-enrichment} corresponds to the identity-on-objects functor from
$\Fel(\cC)$ to the category whose objects are cats and whose morphisms are
functions between leg sets, that is, to the leg functor. No information beyond
$\legs$ is contained in the self-enrichment.
\end{remark}

\section{Additive enrichment}\label{sec:additive}

In this section $\cC$ is a \emph{preadditive} category, that is, a category
enriched in abelian groups, possessing an initial object.

\begin{lemma}\label{lem:Ab-zero}
In a preadditive category every initial object is a zero object. In
particular head and tail of every $\cC$-cat are zero objects, and $\cC$-cats
exist.
\end{lemma}

\begin{proof}
Let $0$ be initial. Then $\cC(0,0)$ is the trivial group, so $\id_0$ is the zero
morphism. For any $f\colon X\to0$, bilinearity of composition gives
$f=\id_0\circ f=0\circ f=0$. Hence $\cC(X,0)$ is a singleton.
\end{proof}

In particular $\cC$ is pointed, so by \cref{rem:pointed-no-chloe} the only
$\cC$-cats satisfying the Chloe Condition are the legless ones, and by
\cref{ex:gideon}(iv) only the zero objects of $\cC$ admit Gideon
Decompositions.


\begin{definition}\label{def:Fel-Ab}
The \emph{additively enriched category of cats} $\Fel_{\Ab}(\cC)$ has as
objects the $\cC$-cats and as hom-groups
\[
\Fel_\Ab(\cC)(K,K')\defeq\prod_{\alpha\in A}\ \bigoplus_{\beta\in B}\cC(L_\alpha,L'_\beta),
\]
whose elements we write as matrices $(f_{\beta\alpha})$ with
$f_{\beta\alpha}\colon L_\alpha\to L'_\beta$ such that for each $\alpha$ only
finitely many $f_{\beta\alpha}$ are nonzero (\emph{column-finite matrices}).
Composition is matrix multiplication, $(gf)_{\gamma\alpha}=\sum_\beta g_{\gamma\beta}f_{\beta\alpha}$,
and identities are identity matrices.
\end{definition}

\begin{lemma}\label{lem:Fel-Ab-category}
$\Fel_\Ab(\cC)$ is a preadditive category.
\end{lemma}

\begin{proof}
The sum defining $(gf)_{\gamma\alpha}$ is finite because $f$ is column-finite.
For fixed $\alpha$, $(gf)_{\gamma\alpha}\neq0$ only if $g_{\gamma\beta}\neq0$ for
one of the finitely many $\beta$ with $f_{\beta\alpha}\neq0$, and each of these
columns of $g$ is finite; hence $gf$ is column-finite. Associativity,
unitality and bilinearity follow from the corresponding properties in $\cC$
and the finiteness of all sums involved.
\end{proof}

\begin{theorem}\label{thm:Ab-collapse}
Let $\cC$ be preadditive with an initial object $0$.
\begin{enumerate}
\item In $\Fel_\Ab(\cC)$, every cat $K$ is the coproduct of the unipedal cats
  $\iota(L_\alpha)$, the $\alpha$-th coprojection being the matrix whose only
  possibly nonzero entry is $\id_{L_\alpha}$, in row $\alpha$; if $K$ is a kitten,
  this is a biproduct.
\item In $\Fel_\Ab(\cC)$ the cephalic cat $\Cephalic=\iota(0)$ is isomorphic to
  the legless cat $\Legless$.
\item If $\cC$ has binary biproducts, then $\iota(X\oplus Y)\iso\langle X,Y\rangle$ in
  $\Fel_\Ab(\cC)$.
\item Consequently, the number of legs is not an isomorphism invariant in
  $\Fel_\Ab(\cC)$, and there is no functor $P$ from the underlying category of
  $\Fel_\Ab(\cC)$ to $\Set$ with $P(K)$ in bijection with $\legs(K)$ for all
  $K$.
\item The underlying ordinary category of $\Fel_\Ab(\cC)$ is not equivalent to
  $\Fel(\cC_0)$, where $\cC_0$ is the underlying ordinary category of $\cC$.
\end{enumerate}
\end{theorem}

\begin{proof}
(i) By definition,
\[
\Fel_\Ab(\cC)(K,M)=\prod_\alpha\bigoplus_\beta\cC(L_\alpha,M_\beta)=\prod_\alpha\Fel_\Ab(\cC)(\iota L_\alpha,M),
\]
and the bijection is precomposition with the stated coprojections. For finite
$A$ a coproduct in a preadditive category is a biproduct.

(ii) $\Fel_\Ab(\cC)(\Cephalic,\Legless)=\bigoplus_{\varnothing}=0$ and
$\Fel_\Ab(\cC)(\Legless,\Cephalic)=\prod_\varnothing=0$, and
$\Fel_\Ab(\cC)(\Cephalic,\Cephalic)=\cC(0,0)=0$; so the zero morphisms are mutually
inverse.

(iii) The matrices $\binom{p_X}{p_Y}\colon\iota(X\oplus Y)\to\langle X,Y\rangle$
and $(i_X\ i_Y)\colon\langle X,Y\rangle\to\iota(X\oplus Y)$, formed from the
biproduct projections and injections, are mutually inverse by the biproduct
identities.

(iv) By (ii), $\Cephalic\iso\Legless$ although $|\legs\Cephalic|=1\neq0=|\legs\Legless|$.
A functor preserves isomorphisms, so $P(\Cephalic)\iso P(\Legless)$, which is
incompatible with $P(\Cephalic)\iso1$ and $P(\Legless)\iso\varnothing$.

(v) The underlying category of $\Fel_\Ab(\cC)$ is preadditive and has the
initial object $\Legless$ (by (i), as the empty coproduct), which is therefore a
zero object by \cref{lem:Ab-zero}. In $\Fel(\cC_0)$ the initial object
$\Legless$ is strict (\cref{prop:legless-initial}); if it were also terminal,
then every cat would admit a morphism to $\Legless$ and hence be legless,
which fails for $\iota(0)$. So $\Fel(\cC_0)$ is not pointed, and an equivalence
would preserve pointedness. (For instance, if $\cC$ is the trivial preadditive
category, $\Fel_\Ab(\cC)$ is equivalent to the terminal category, every cat
being a coproduct of zero objects, while $\Fel(\cC_0)\iso\Set$.)
\end{proof}

\begin{example}\label{ex:free-abelian}
Let $\cC$ be the full subcategory of $\Ab$ on the objects $0$ and $\mathbb Z$. The
functor sending a cat $K$ to $\bigoplus_\alpha L_\alpha$ is an equivalence
between $\Fel_\Ab(\cC)$ and the category of free abelian groups. Indeed, it is
fully faithful because a homomorphism
$\bigoplus_{\alpha}L_\alpha\to\bigoplus_\beta L'_\beta$ is determined by its
restrictions to the summands, each of which lands in a finite sub-sum; and it
is essentially surjective since every free abelian group is of the form
$\bigoplus_A\mathbb Z$. The legs equal to $0$ are invisible, in accordance with
\cref{thm:Ab-collapse}(ii).
\end{example}

\begin{remark}\label{rem:enrichment-moral}
\Cref{thm:Ab-collapse} shows that the naive enrichment of cat theory over
abelian groups does not merely enrich the theory of \cref{ch:cats} but changes
it: legs may be merged and split, and a leg equal to the head may be
amputated without changing the isomorphism class of the cat. The
felinization of an enriched category is therefore not the enrichment of the
felinization. We do not develop enrichment over a general base $\cV$. One
would expect the formula $\prod_\alpha\coprod_\beta\cC(L_\alpha,L'_\beta)$ for
hom-objects to define a $\cV$-category when $\cV$ has small products and small
coproducts and $\otimes$ preserves coproducts in each variable, but we have not
verified the axioms and make no claim to this effect.
\end{remark}
